\section{Introduction}
\label{sec:intro}

The Lightning Network is a payment network built on Bitcoin. Its nodes hold their operators' funds in two-party
payment channels, each backed by one output on the Bitcoin chain. When both parties agree, a channel closes
cooperatively with one ordinary transaction. When they cannot, because a peer is offline, a payment is stuck or the two
implementations disagree, either party can force-close the channel: it broadcasts the latest commitment transaction,
signed in advance by both, and then claims its share through further pre-signed and time-locked
transactions~\cite{poon2016bitcoin,bolt3,bolt5}. The fees and delays of those transactions fall on the operators, and
a transaction that confirms too late can lose a pending payment or let a cheating peer
win~\cite{harris2020flood,riard2020timedilation,riard2023replacement}. The fee design of these transactions is
changing now. Bitcoin Core 31.0 (April~2026) removed the CPFP carve-out, the mempool rule that anchor channels were
designed around~\cite{core31}, and in May~2026 the Lightning specification adopted zero-fee commitments, which rely on
newer relay rules instead~\cite{bolts1228zerofee,bip431,core28}. This paper measures the old design, across the fee
regimes it lived through, before the migration, as a baseline for judging the new one.

Force closes have been counted~\cite{bitmex2020closures,zabka2022empirical,grotschla2025lifecycle} and
predicted~\cite{antonelli2026closures}, but what they cost has not been measured across fee regimes. The questions
that shaped the design are argued in developer forums from a few nodes' logs or from none: whether owners use their
anchor outputs to bump fees, whether third parties sweep them at a profit, and whether fee spikes force channels
closed~\cite{daftuar2024v3,lnd8997}. Our previous paper measured the cost of force closes over a few months of low
fees in 2026~\cite{kumar2026clientdiversity}.

We measure it from the chain.
We follow every close of every public channel since Lightning's first channels in January~2018, find the force
closes of channels outside the public graph by the fingerprint of their commitment transaction, and follow each of
their \NOutputs{} outputs to the transaction that spent it: \NClosesAll{} closes in all, \NForceAll{} of them force
closes. Fee statistics for every block, from our own node, place each transaction in the fee market of its block.
The main answer is that fee spikes do not bring a wave of force closes. In the week of a spike and the week after,
force closes were \PlEvAllNow{} times their level in the four weeks before, while ordinary weeks of the same years
show \PlEvAllNowMed. That holds for pre-anchor channels, whose commitment fee is fixed when it is signed, as for anchor
channels, whose fee can be raised at broadcast. The cost of a force close lies elsewhere: in its ratio to a cooperative close, in
the days the closer's funds stay locked, and in the exposure of pending payments when fees are high. The wait
outweighs the fees only in low-fee years, and the exposure has rarely become a loss. We make five contributions.

\begin{itemize}
\item \textbf{A census with rates.} Force closes became rarer per channel, and private channels now produce most of
them. Of the public channels funded in 2019, \CohortOneYearNineteen\% were force-closed within a year; of those
funded in 2025, \CohortOneYearTwentyFive\%. The force share of public closes fell from \ForceShareTwentyThree\% in
2023 to \ForceShareTwentySix\% in 2026. In 2026 force closes of channels outside the public graph outnumbered public
ones \UnlOverPubTwentySix{} to one (\S\ref{sec:census}).
\item \textbf{The cost of a force close in fees, time and block space.} A force close costs
\CostRatioMin--\CostRatioMax{} times a cooperative close in fees and, since 2022, about three times its block space; the
closer waits a median of \ToLocalMedDays{} days for its own balance, which outweighs the fees only in the low-fee
years 2025--26. Counting the open as well, a public channel's whole life cost a median of \LcMed{}~sat on-chain. When
fees are high, HTLC timeouts confirm late, and anchor channels' timeouts less often than pre-anchor ones. Yet
forwarders lost few payments: \SplitResForwarderPayments{} since \SplitResFromYear, all with the preimage on-chain
first. The only losses in the order a replacement-cycling attack leaves were a test that one operator ran against its
own node, and 2026's one wave of penalties matches a splicing bug already fixed in Core Lightning
(\S\ref{sec:cost}).
\item \textbf{The anchor money trail.} Anchor outputs became third-party dust. Owners spent
\OwnerPctTwentyTwo\% of anchors in 2022 and \OwnerPctTwentySix\% in 2026; third parties sweep the rest, from 2025 at
the first block they can, and fees took \MinersShareTwentySix\% of the swept value in 2026. No P2WSH anchor was left
unspent, while no third party has swept a simple-taproot anchor, so the premise of an open LND proposal to stop
sweeping anchors holds only for the older kind (\S\ref{sec:anchors}).
\item \textbf{Fee regimes.} Fee spikes did not bring a wave of force closes, judged against placebo weeks of
the same years. Fixed commitment fees overpaid in calm
markets, and all of Lightning's closes and required spends took at most
\BillFeePctMax\% of Bitcoin's transaction fees in any year (\S\ref{sec:fees}).
\item \textbf{The first zero-fee commitments.} The first zero-fee commitment closed on \ZfFirst, before the
specification was merged. Across the anchor force closes we observe, the new format would have saved
\CfSavingPct\% of the anchor mechanism's block space (\S\ref{sec:zerofee}).
\end{itemize}

This is the third paper of a programme on the Lightning Network's resilience. The first asked whether client diversity
protects the network when one implementation fails~\cite{kumar2026clientdiversity}, the second who patches
it~\cite{kumar2026upgradedynamics}, and this one what a failed channel costs on-chain. All data behind the figures,
tables and in-text numbers are released with the code. We report aggregates only: no node, address or transaction is
named.
